\documentclass[conference,a4paper]{IEEEtran}
\usepackage[utf8]{inputenc}
\usepackage[T1]{fontenc}
\usepackage[spanish,english]{babel}
\usepackage{graphicx}
\usepackage{booktabs,multirow}
\usepackage{siunitx}
\usepackage[protrusion=true,expansion=true,final]{microtype}
\usepackage{hyperref}
\usepackage{flushend}
\usepackage{listings}
\usepackage{xcolor}
\usepackage[a4paper,top=2.5cm,left=2.5cm,right=2.0cm,bottom=2.0cm,showframe=false]{geometry}

\definecolor{lightgray}{rgb}{0.95,0.95,0.95}
\definecolor{mediumgray}{rgb}{0.85,0.85,0.85}

\lstset{
  basicstyle=\ttfamily\footnotesize,
  breaklines=true,
  captionpos=b,
  commentstyle=\color{gray},
  escapeinside={\%*}{*)},
  keywordstyle=\color{blue},
  stringstyle=\color{purple},
  numbers=left,
  numberstyle=\tiny\color{gray},
  numbersep=5pt,
  tabsize=2,
  frame=single,
  framesep=5pt,
  xleftmargin=15pt,
  backgroundcolor=\color{lightgray},
  rulecolor=\color{mediumgray},
  showstringspaces=false
}

\begin{document}
\selectlanguage{spanish}
\title{Documento Test 6}
\author{Test Author}
\maketitle

\begin{abstract}
Prueba con configuración de listings corregida.
\end{abstract}

\section{Introducción}
Probando valores numéricos: \SI{16}{\giga\byte}, \SI{1.7}{\milli\second}.

\begin{lstlisting}[language=Python]
# Un ejemplo de código Python
def hello_world():
    print("Hola mundo")
\end{lstlisting}

\end{document}
EOF < /dev/null
